% Chapter 11: DeepONet: Engineering Branch-Trunk Networks in PyTorch
% Part II of the book.
% STATUS: book-standard prose (Part II).

\chapter{DeepONet: Engineering Branch-Trunk Networks in PyTorch}
\label{ch:part2-deeponet-engineering-branch-trunk-networks}
\label{sec:deeponet_impl}
% -----------------------------------------------------------------------------

DeepONet (Deep Operator Network) \cite{lu2021learning} turns the existence statement of
\chapref{ch:part1-operator-approximation-theory-and-universa} into a module one can import.
Theorem~\ref{thm:chen_chen} of \cite{chen1995universal} established that a continuous operator
on a compact set of a function space is approximable to accuracy $\varepsilon$ by a finite sum
of separable terms \eqref{eq:chen_chen_form}: $m$ readings of the input, $p$ functions of the
output coordinate, and $n$ nonlinear mixings between them. Equation \eqref{eq:ch04-deeponet}
replaces the mixings by one trainable network and the output functions by another, joined by a
single inner product. Between those two displays sits this chapter's problem. Four choices
precede the first gradient step: the sensor count $m$ and its locations, the shared latent
width $p$, the depth and hidden width of each sub-network, and the batched layout in which the
contraction runs on a graphics processing unit (GPU). The first three price one term apiece in
the bound of Theorem~\ref{thm:deeponet_error}---encoding
$\mathcal{E}_{\text{sensor}}(m)$, reconstruction $\mathcal{E}_{\text{basis}}(p)$, network
$\mathcal{E}_{\text{DNN}}(\theta)$ in \eqref{eq:ch04-error-split}---and the fourth buys speed
rather than accuracy. The first three are also irreversible: $m$ is welded into one parameter
tensor, $p$ into two, and the hidden widths into all of them, so branch-trunk design is a
budget allocation settled before the first epoch rather than a hyperparameter sweep. Only the
fourth choice stays open afterwards, since a layout can be rewritten without touching a single
weight\index{DeepONet!branch-trunk implementation
in PyTorch}.

The reader should bring PyTorch \cite{paszke2019pytorch} at the level of writing an
\texttt{nn.Module} of \texttt{nn.Sequential} stacks, plus the GPU vocabulary of device
residency and batched matrix multiplication; the layout, loss, and normalization conventions of
\chapref{ch:part2-the-practitioner-s-paradigm-infinite-dimen} are used without restatement,
including the relative $L^{2}$ loss \eqref{eq:rel_l2_loss}, the metric that prices
$\mathcal{E}_{\text{DNN}}(\theta)$ in an experiment. Two Part~I conditions make the architecture
well posed at all: point evaluation must be a bounded functional on the input space, the
Sobolev embedding condition of Remark~\ref{rmk:ch8-inclusions}, and the input ensemble must be
efficiently encodable, which is where the rate $\mathcal{E}_{\text{sensor}}(m)\lesssim m^{-s/d}$
of \eqref{eq:ch8-sensor-rate} comes from. Panel (b) of Figure~\ref{fig:architectures} shows the
shape of the design: fixed input measurements, free output queries, one inner product joining
them.

The route is deliberately asymmetric. Section~\ref{sec:ch22-blueprint} fixes the functional
form, walks the tensor shapes through the listing, and does the parameter and multiply-add
bookkeeping that ties the module to the DeepONet line of the cost ledger
\eqref{eq:ch8-costs}. Section~\ref{sec:ch22-sensors} treats the sensor set as a conditioning
problem rather than a checklist, comparing three node families along one axis: what each
protects the encoding against. Section~\ref{sec:ch22-latent} reads $p$ as a rank constraint,
states what a stationary error floor means for the three terms, and records how the module
fails.

% -----------------------------------------------------------------------------
\section{Architectural Blueprint}
\label{sec:ch22-blueprint}
% -----------------------------------------------------------------------------

Let the input function $u$ be sampled at $m$ sensor locations
$\{x_1, \dots, x_m\} \subset \Omega$, collected as
$\mathbf{u}_{\mathrm{sensor}} = \big(u(x_1), \dots, u(x_m)\big) \in \mathbb{R}^m$. The branch
network maps that vector to a latent code $\mathbf{b}(u) \in \mathbb{R}^p$, where $p$ is the
shared \index{latent width@latent width $p$}latent width; simultaneously the trunk network
evaluates a query coordinate $y \in \Omega \subset \mathbb{R}^d$ and returns
$\mathbf{t}(y) \in \mathbb{R}^p$. The prediction at $y$ is the pairing of the two codes plus a
scalar offset $b_0 \in \mathbb{R}$:
\begin{equation}
  \label{eq:ch22-deeponet-contract}
  \mathcal{G}_\theta(u)(y) = \sum_{k=1}^p b_k(u) \, t_k(y) + b_0.
\end{equation}
Equation \eqref{eq:ch22-deeponet-contract} is \eqref{eq:ch04-deeponet} restated for
implementation, and the summation index carries the whole design constraint. Because $k$ is
summed against both factors, the branch and the trunk must emit the same number of components;
the module cannot express a mismatch, since both stacks terminate in a linear layer whose
output width is the single constructor argument \texttt{latent\_p}. The asymmetry Part~I
stressed appears one level down, in what the two constructor arguments mean: \texttt{coord\_dim}
fixes only the dimension of a space the user may query anywhere in, whereas \texttt{m\_sensors}
fixes the length and the ordering of a vector whose entries are prescribed by the data
pipeline\index{DeepONet!trunk network}\index{DeepONet!branch network}. The output side is free;
the input side is not.

Evaluated on $P$ query points at once, the sum of outer products becomes one matrix product
\index{batched matrix multiplication}
\begin{equation}
  \label{eq:ch22-batch-contraction}
  \mathbf{G} = \mathbf{B} \mathbf{T}^T + b_0 \in \mathbb{R}^{B \times P}, \qquad \mathbf{B} \in \mathbb{R}^{B \times p}, \quad \mathbf{T} \in \mathbb{R}^{P \times p},
\end{equation}
where $B$ is the batch dimension, that is, the number of independent input functions carried
in one step, $\mathbf{B}$ stacks their branch codes, and $\mathbf{T}$ stacks the trunk
evaluations at the query points. The rank structure in \eqref{eq:ch22-batch-contraction} is
the reason the architecture is cheap: one branch pass serves every query point of a sample,
and when the query set is shared across the batch the trunk is evaluated once per step rather
than $B$ times.

Listing~\ref{lst:ch22-deeponet} realizes both halves as multilayer perceptrons (MLPs) of depth
three and hidden width $h$ (\texttt{hidden\_dim}, default 128), separated by the Gaussian
error linear unit (GELU), with the trainable offset \texttt{bias} added to the $(B,P)$ output
in the last line. Two choices in the listing deserve a reason. The activation is a regularity
decision before it is a training decision: the discretization-error analysis of
\chapref{ch:part1-discretization-invariance-and-continuum-li} caps the achievable decay at the
regularity of the \emph{least} smooth object in the pipeline. The rectified linear unit (ReLU)
has regularity only just below $H^{3/2}$, so with ReLU in place of GELU the observed decay
saturates once the inputs are smoother than that, whereas a smooth activation imposes no such
ceiling and preserves the rate $N^{-s}$ of \eqref{eq:ch06-rate}. The qualification matters:
the ceiling binds only when the activation is the roughest object in the path, and a rough
input field or a non-periodic coordinate ramp can sit below it. The additive scalar bias,
initialized at zero, absorbs the ensemble mean of the output field so that the latent code is
spent on variation rather than on offset; the pointwise normalizer of
\chapref{ch:part2-the-practitioner-s-paradigm-infinite-dimen} does the same job on the input
side, and together they determine what is left for $\mathbf{b}(u)$ and $\mathbf{t}(y)$ to
represent---the part of the field that actually varies.

\begin{lstlisting}[language=Python, caption={\textsf{DeepONet} in PyTorch: two independent stacks joined only at the latent width $p$. The constructor fixes the sensor count $m$, the coordinate dimension $d$, and $p$; the closing \texttt{einsum} (Einstein summation) calls are meant to carry out the batched contraction of \eqref{eq:ch22-batch-contraction}, and the discussion after the listing explains why, as printed, they do not.}, label={lst:ch22-deeponet}]
import torch
import torch.nn as nn

class DeepONet(nn.Module):
    """Deep Operator Network with unstacked branch and trunk MLPs."""
    def __init__(self, m_sensors: int, coord_dim: int, latent_p: int, hidden_dim: int = 128):
        super().__init__()
        self.latent_p = latent_p
        
        # Branch network: processes input function evaluated at m fixed sensors
        self.branch = nn.Sequential(
            nn.Linear(m_sensors, hidden_dim),
            nn.GELU(),
            nn.Linear(hidden_dim, hidden_dim),
            nn.GELU(),
            nn.Linear(hidden_dim, latent_p)
        )
        
        # Trunk network: processes continuous query coordinates y in R^d
        self.trunk = nn.Sequential(
            nn.Linear(coord_dim, hidden_dim),
            nn.GELU(),
            nn.Linear(hidden_dim, hidden_dim),
            nn.GELU(),
            nn.Linear(hidden_dim, latent_p)
        )
        
        # Trainable output bias scalar
        self.bias = nn.Parameter(torch.zeros(1))

    def forward(self, u_sensors: torch.Tensor, y_coords: torch.Tensor) -> torch.Tensor:
        """
        u_sensors: (B, m) - evaluations at m sensor points
        y_coords:  (B, P, coord_dim) or (P, coord_dim) - evaluation points
        Returns:   (B, P) - predicted field values at query locations
        """
        B_features = self.branch(u_sensors) # Shape: (B, p)
        
        if y_coords.dim() == 2:
            # Query coordinates shared across entire batch: (P, coord_dim)
            T_features = self.trunk(y_coords) # Shape: (P, p)
            out = torch.einsum('bp, pp -> bp', B_features, T_features) # (B, P)
        else:
            # Query coordinates vary per sample: (B, P, coord_dim)
            T_features = self.trunk(y_coords) # Shape: (B, P, p)
            out = torch.einsum('bp, bpp -> bp', B_features, T_features)
            
        return out + self.bias
\end{lstlisting}

What the reader should take from Listing~\ref{lst:ch22-deeponet} is not the module skeleton but
the two contraction calls, because they are where tensor-shape reasoning is checked. The
\texttt{einsum} operation implements the Einstein summation convention: every axis of every
operand receives a letter, axes sharing a letter are matched, and letters absent from the
output label are summed\index{Einstein summation convention (einsum)}.
In the first path, query coordinates shared across the batch arrive as $(P, d)$ and produce
trunk features of shape $(P, p)$, so the intended operation is exactly
\eqref{eq:ch22-batch-contraction}; in the second, per-sample coordinates $(B, P, d)$ produce
$(B, P, p)$ and the batch label must be matched as well. Note, however, that both strings as
printed reuse the letter \texttt{p} for the latent axis and for the query axis, and a letter
repeated inside a single operand takes the diagonal of those two axes and forces them to the
same length. The consequence is stated plainly: the contraction of
\eqref{eq:ch22-batch-contraction} is obtained only if the query and latent axes are labeled
distinctly---\texttt{bp,qp->bq} for a shared grid and \texttt{bp,bqp->bq} for per-sample
coordinates---and a caller who asks for $P = 256$ queries with $p = 128$ will meet a shape
error rather than a silently wrong field, while a caller who happens to take $P = p$ receives
numbers, wrong ones, with no complaint. The listing is reproduced as published; the
reconciliation is recorded here because no unit test catches this failure by
itself\index{tensor layout!branch-trunk contraction}.

The arithmetic of the two stacks follows directly from their shapes. Each linear layer of
width pair $(r, c)$ stores $rc$ weights and $c$ biases, so the branch contributes
$mh + h + h^{2} + h + hp + p$ and the trunk $dh + h + h^{2} + h + hp + p$, giving
\begin{equation}
  \label{eq:ch22-params}
  N_{\mathrm{par}} = 2h^{2} + h\,(m + d + 2p + 4) + 2p + 1 .
\end{equation}
With the listing's default $h = 128$ and a representative $d = 2$, $p = 64$ and $m = 100$
sensors, that is
$32{,}768 + 128 \cdot 234 + 129 \approx 6.3 \times 10^{4}$ parameters, of which the two hidden
square layers supply just over half. Two readings of \eqref{eq:ch22-params} matter in practice.
The sensor budget $m$ enters linearly and only through the first branch layer, so doubling the
readings adds $hm = 1.28 \times 10^{4}$ entries---small next to the cost of the solver runs
that produced the training pairs, which is the comparison Part~I's cost bookkeeping
\eqref{eq:ch01-cost-bookkeeping} insists on. The latent width $p$ enters as $2hp$, that is, once per
stack, so widening the shared code is also cheap; what widening $p$ does \emph{not} buy is
discussed in Section~\ref{sec:ch22-latent}.

Per step, the multiply-add count for a batch of $B$ samples evaluated at $P$ query points is
\begin{equation}
  \label{eq:ch22-query-cost}
  \Phi(B, P) = B \Big( \underbrace{mh + h^{2} + hp}_{\text{branch}} + P \underbrace{(dh + h^{2} + hp)}_{\text{trunk, per query point}} + Pp \Big),
\end{equation}
the last term being the entries of the rank-$p$ contraction. The count is the per-sample one:
the trunk work is charged $B$ times. On a query grid shared across the batch, drop the factor
$B$ from the two trunk terms only---the contraction stays $B P p$---and the single trunk pass of
\eqref{eq:ch22-batch-contraction} is what the ledger should charge. Setting $d_b = d_t = h$,
\eqref{eq:ch22-query-cost} is
the DeepONet row of the cost ledger \eqref{eq:ch8-costs}, $\Theta(m d_b + p d_t + p N_{\mathrm{out}})$,
spelled out: the branch pass does not grow with the number of output
points, and the output pass grows linearly in it. Memory behaves differently\index{memory
footprint!trunk activations}: the trunk activations dominate, $P p$ entries for a shared query
grid and $B P p$ for per-sample grids---so $65{,}536$ query points on a $256 \times 256$ output
grid with $p = 64$ reserve about $4.2$ million floats, roughly $17$ megabytes per sample before
autograd storage.
The standard remedy is to split the query set into chunks over $P$, since
\eqref{eq:ch22-batch-contraction} makes the output entries independent of one another; the
branch code, once computed, is reused across all chunks.

% -----------------------------------------------------------------------------
\section{Sensor Placement Best Practices}
\label{sec:ch22-sensors}
% -----------------------------------------------------------------------------

The accuracy of the branch network depends on the sensor configuration
$\{x_j\}_{j=1}^{m}$, and the dependence is not a detail of data preparation but a statement
about the architecture: the sensor set is the part of a DeepONet that the network never sees.
The first linear layer consumes $m$ numbers in a fixed order, so the locations enter only
through their count. What the branch learns is the statistics of the readings, and if two
functions in the input ensemble agree at $\{x_j\}_{j=1}^{m}$ while differing elsewhere, the
model is being asked to distinguish inputs that its own argument list says are identical. The
encoding term $\mathcal{E}_{\text{sensor}}(m)$ of \eqref{eq:ch04-error-split} measures exactly
that ambiguity, and \eqref{eq:ch8-sensor-rate} gives its size for ensembles bounded in
$H^{s}(\Omega)$ with $s > d/2$: the error decays like $m^{-s/d}$, so halving it costs a
factor $2^{d/s}$ in sensors---four times the readings for a two-dimensional ensemble with one
Sobolev derivative, and about $1.6$ times for the same two-dimensional ensemble with three
derivatives\index{sensor placement!cost of halving the
encoding error}\index{point evaluation!conditioning of the encoding}.
Regularity, not resolution, sets the price, and the same inequality $s > d/2$ that makes
$\mathcal{E}_{\text{sensor}}$ finite is the boundedness condition of
Remark~\ref{rmk:ch8-inclusions}: for rougher inputs, point evaluation is not a continuous
functional, sensor noise is amplified without control, and no placement strategy rescues the
architecture.

Three families of node sets are used in practice, and they differ along one axis---how node
density is distributed over $\Omega$, and correspondingly which failure each guards
against\index{sensor placement!equispaced, space-filling, and clustered node sets}.
\emph{Equispaced grids} put the density where the geometry is simplest and are the natural
choice for smooth, periodic fields: the sensor vector is precisely what a discrete Fourier
transform consumes, so a band-limited ensemble is recovered from its samples and the encoding
error is the tail beyond the sensor Nyquist rate\index{Nyquist rate!sensor sampling}. The
exposure is equally clear. Once the ensemble carries energy above that rate, the readings do
not merely omit it---the sampled representation folds it back onto the resolved coefficients,
so the branch is trained on a corrupted encoding rather than an incomplete one, and neither a
wider trunk nor a longer run repairs it\index{aliasing!sensor sampling}. Steep gradients in
shock-forming data are the usual casualty.

\emph{Latin hypercube sampling} (LHS) inverts the priority: it maximizes space-filling coverage
across irregular or multi-dimensional domains, where no lattice is natural\index{Latin hypercube
sampling (LHS)}. The construction divides each coordinate axis into $m$ equal-probability
strata, places exactly one sensor in every stratum of every axis, and pairs the resulting
one-dimensional samples by a random permutation, so no two sensors share a slab in any
coordinate. The property bought is isotropy: the worst-case gap between nodes, the fill
distance of \chapref{ch:part1-discretization-invariance-and-continuum-li}---its mesh parameter
$h_N$ evaluated on a node set of $m$ sensors rather than on a mesh of $N$ grid
points\index{fill
distance!sensor sets}---is controlled in every direction at once rather than only along lattice
lines, and there is no preferred direction along which aliasing can organize itself. The
property paid for is the logarithm in the coverage asymptotics---the fill distance of $m$
random points in $d$ dimensions is of order
$(\log m / m)^{1/d}$ where a grid achieves $m^{-1/d}$---so a random set needs somewhat more
sensors for the same worst-case gap. That cost is often worth accepting, because the error
analysis of \cite{lanthaler2022error} prices ensembles whose sensors are drawn at random from
a measure rather than tuned by hand, so the probabilistic construction sits inside the theorem
rather than outside it.

\emph{Chebyshev extrema nodes} cluster the density at the boundary, where non-periodic problems
put their structure:
\begin{equation}
  \label{eq:ch22-chebyshev-nodes}
  x_j = \frac{1}{2} \left( 1 - \cos\left( \frac{(j-1)\pi}{m-1} \right) \right), \quad j = 1, \dots, m.
\end{equation}
Equation \eqref{eq:ch22-chebyshev-nodes} maps the extrema of the Chebyshev polynomial onto
$[0,1]$; consecutive spacing is of order $1/m$ in the interior but shrinks to order
$1/m^{2}$ at the endpoints, and the set contains both endpoints, so the branch reads the trace
values that Dirichlet data prescribe\index{Chebyshev extrema nodes}. The mechanism at work is
classical polynomial interpolation, and the branch inherits its conditioning: the encoding is
stable only if the function reconstructed from the readings stays close to the true one in the
gaps between sensors. Reconstruction from equispaced samples carries a Lebesgue constant that
grows exponentially with $m$, and the interpolants of smooth non-periodic functions oscillate
with growing amplitude near the boundary---Runge's phenomenon\index{Runge phenomenon}, whose
mitigation is the property this node set is selected for. For the extrema nodes the nodal
polynomial $\prod_j (x - x_j)$ remains small and nearly equioscillating on the interval, so the
growth of that constant is logarithmic instead. That is precisely the well-posedness the
encoding term assumes: an ill-conditioned sampling set makes $\mathcal{E}_{\text{sensor}}(m)$
large and, worse, non-uniform, concentrating error where the sensor gap is widest.

The three families are therefore not ranked. Equispaced nodes minimize the worst-case gap along
the coordinate lines and buy it back as directional aliasing; random space-filling nodes pay a
logarithm in gap size and buy isotropy plus a theory that assumes randomness; clustered nodes
trade interior spacing for boundary resolution and buy stability of the reconstruction in
non-periodic problems. The choice is made by the input ensemble: its periodicity, its Sobolev
regularity, and whether the operator of interest is sensitive to traces. Two practical
constraints close the section. Because $m$ is welded into a parameter tensor, re-placing the
sensors invalidates the trained branch and the model must be retrained, not patched. And
because the module takes the readings as an ordered vector, permuting the sensor axis at
inference silently changes the learned operator; nothing in the shapes complains, so the
convention belongs in the checkpoint alongside the weights.

% -----------------------------------------------------------------------------
\section{Latent Width, Output Structure, and Failure Modes}
\label{sec:ch22-latent}
% -----------------------------------------------------------------------------

The parameter $p$ is a rank, not a width, and reading it as a width obscures what it
actually constrains. The right-hand identity of \eqref{eq:ch8-deeponet-kernel} makes the
statement precise: the kernel realized by \eqref{eq:ch22-deeponet-contract} has rank at most
$p$ in the pair of variables, and for every input the predicted field lies in the
$p$-dimensional span of the trunk functions plus a constant. The reconstruction term
$\mathcal{E}_{\text{basis}}(p)$ of \eqref{eq:ch04-error-split} is therefore a subspace
approximation problem, and the mean-square-optimal $p$-dimensional subspace for the output
ensemble is spanned by the leading eigenvectors of its covariance operator---the
Eckart--Young argument in finite dimension, its Karhunen--Lo\`eve counterpart for
function-valued ensembles, which is exactly the spectral decay through which
\cite{lanthaler2022error} prices this term\index{trunk basis!optimal rank-$p$ subspace}
\index{covariance operator!spectral decay}. The engineering use of the fact is a diagnostic
available before training: compute the fraction of output energy captured by the leading $p$
empirical modes of the training targets. If that fraction does not approach one, the ensemble
is not $p$-dimensionally encodable and enlarging $p$ is the only remedy---more epochs cannot
reduce a term that \eqref{eq:ch04-error-split} attributes to the basis. The same logic gives
the triage for a stalled run, and which curve is stalled carries the information. A basis
bottleneck binds the \emph{training} loss, because the model class of
\eqref{eq:ch22-batch-contraction} produces a matrix of rank at most $p$: the curve settles at a
level that more epochs cannot lower. A shortfall in the readings usually shows the other way
round, with the training loss still falling while the validation error floors, since a fixed
sensor set can be partly memorized on the training ensemble and cannot be recovered on a fresh
sample. Each hypothesis is then tested by the intervention that moves only it, at fixed data:
raise $p$ with $m$ fixed, raise $m$ with $p$ fixed, or raise the training budget and the
capacity with the architecture fixed. Only the last touches $\mathcal{E}_{\text{DNN}}(\theta)$.

The listing implements the unstacked form named in its own docstring, in which one latent code
and one trunk serve the whole output; the return shape $(B,P)$ is a scalar output field, and a
field with $d_u > 1$ output channels---$d_u$ is Part~I's channel count for
$\mathcal{U}$---either enters through additional query coordinates or through the stacked
variant of \cite{lu2021learning}, in which each output channel carries its own branch and
trunk pair\index{DeepONet!stacked and unstacked variants}. The comparison is parameter count
against sharing: stacking multiplies the branch cost by $d_u$ and lets each channel have its
own basis, while the unstacked form ties all channels to the same $p$ trunk functions and
keeps the contraction of \eqref{eq:ch22-batch-contraction} intact. Channels that share spatial
structure---the pressure and velocity components of one flow, say---can be read off a single
shared code and pay one branch pass; channels whose spatial scales differ by an order of
magnitude are then forced through the same $p$ basis functions, and the constraint surfaces as
a floor on the harder channel rather than as an error message. Which of the two regimes an
application is in is a property of the output ensemble, measured the same way as in the
diagnostic above, not a choice the architecture makes for the user.

The remaining failure modes are worth collecting, because each is silent or nearly so. A trunk
whose outputs are nearly linearly dependent has effective rank below $p$: the weights are
being spent on directions the branch cannot reach, visible as rapidly decaying singular values
of the matrix $\mathbf{T}$ in \eqref{eq:ch22-batch-contraction}, and invisible in the loss
curve. Trunk inputs in physical units are another: a coordinate of magnitude $10^{3}$
dominates the pre-activations of the first trunk layer and the effective input becomes one
dimension plus noise, which is why the grid conventions of
\chapref{ch:part2-the-practitioner-s-paradigm-infinite-dimen} keep coordinates in $[0,1]^d$.
And the trunk must resolve the output at the query spacing $\Delta y$: the basis rank bounds
what the model class contains, but the depth and width of the trunk bound what the trunk can
represent within it, so boundary layers and sharp fronts need trunk capacity as well as
$p$. How the family fares against these constraints on a real benchmark is reported in
Figure~\ref{fig:benchmarks}(b) of
\chapref{ch:part2-canonical-pde-benchmarks-and-reference-wor}\index{Darcy flow!branch-trunk
baseline}, where a branch-trunk network with $m = 100$ sensors is plotted on two-dimensional
Darcy flow beside a convolutional baseline and a spectral operator. The conclusion the panel
carries is the ordering: the branch-trunk trace ends highest of the three and is markedly the
noisiest, which is what a floor set by the encoding or the basis term looks like rather than a
floor set by training---though a convergence plot alone cannot say which of the two terms is
responsible.

\bigskip
\noindent Both restrictions this chapter has been circling---the $m$ prescribed readings on
the input, the $p$ trunk functions on the output---descend from one design act, the pairing of
two codes by a single inner product, and everything the architecture gains and loses is a
consequence of that pairing. \index{Fourier neural operator (FNO)!versus DeepONet}The next
chapter changes what the parameters are functions of---wavenumbers instead of sensor
slots---and with it the cost line: where \eqref{eq:ch22-query-cost} charges a fixed sensor set
and a rank-$p$ readout, the Fourier neural operator (FNO) \cite{li2021fourier} reads every grid
point at every layer and pays $\Theta(d_v N \log N + K d_v^{2})$, with the ledger's own symbols
($d_v$ the lifted width, $K$ the retained modes, $N$ the number of grid nodes), the second row
of \eqref{eq:ch8-costs}, the $\log N$ bought by the fast Fourier transform (FFT) and the torus
required by it. \chapref{ch:part2-fourier-neural-operators-fno-from-scratch} builds that
layer from the same conventions used here: the real-input two-dimensional transform
(\texttt{torch.fft.rfft2}), the two corner slices of the half-spectrum that the conjugacy
condition \eqref{eq:ch05-reality} explains, the mode-wise complex contraction, and the lifting
and projection stacks that surround it. The \texttt{einsum} discipline that governs the
contraction in Listing~\ref{lst:ch22-deeponet} returns there as \texttt{"bixy,ioxy->boxy"},
five letters to keep straight---batch, input channel, output channel, and the two mode
axes---and the habit formed here, checking which letters are matched and which are summed, is
the one that makes the spectral layer readable rather than mysterious.
